\subsubsection{视界谱测度的离散化与共动层析：红移权重、洛伦兹窗与酉谱重构接口}\label{subsubsec:xi-horizon-spectral-measure-tomography}
本小节把“视界 Carath\'eodory--Toeplitz 证书”推进到可用于谱重构与层析扫描的精确层面：在 RH 前提下，视界谱测度不仅存在而且可取为纯点原子；其原子权重被 Cayley 紧化的 Jacobian 结构性地强制为 \((1+\gamma^{2})^{-1}\)（红移权重）。
进一步地，引入共动平移参数 \(x_0\) 得到一族洛伦兹窗（Cauchy 核）加权的离散谱测度，它把“视界毛发（多极矩）”组织为可平滑扫描的矩序列族，并与文稿的“拉马努金型对称＝相位化闭包（unitary slice）”制度完全相容。
\par
以下沿用附录 \ref{subsubsec:app-horizon-weyl-carath-toeplitz} 的记号：\(E_{\zeta}(z)=\Xi_{\zeta}(\tfrac12+\mathrm{i}z)\)、\(m_{\zeta}(z)=-E_{\zeta}'(z)/E_{\zeta}(z)\)，Cayley 紧化为 \(w=\mathcal{C}(z)\) 且 \(\mathcal{C}^{-1}:\mathbb{D}\to\mathbb{H}\)。
视界 Carath\'eodory 对数导数函数为（见附录式 \eqref{eq:app-horizon-carath-def}）
\[
  \mathscr{C}_{\zeta}(w):=-\frac{\Xi_{\zeta}'(s(w))}{\Xi_{\zeta}(s(w))},
  \qquad s(w)=\tfrac12+\mathrm{i}\,\mathcal{C}^{-1}(w).
\]

\paragraph{RH 下的视界谱测度原子化（红移权重）}
在 RH 成立时，\(\mathscr{C}_{\zeta}\) 属于 Carath\'eodory 类并具有纯原子谱测度；其闭式原子化正规形已由附录定理 \ref{thm:app-horizon-spectral-measure-atomic} 给出。
为后续层析接口起见，这里将其核心公式以可直接引用的形式提取如下：
存在纯点正测度 \(\mu_{\zeta}\)（谱测度）使得
\begin{equation}\label{eq:xi-horizon-mu-atomic-recall}
\boxed{\;
\mu_{\zeta}
=
\sum_{\gamma\in\Gamma_+}\frac{\mathrm{mult}(\gamma)}{1+\gamma^2}\Bigl(\delta_{\xi_{\gamma}}+\delta_{\xi_{\gamma}^{-1}}\Bigr),
\qquad \xi_{\gamma}:=\mathcal{C}(\gamma)\in\TT,
\;}
\end{equation}
并且其 Fourier 矩（视界多极矩）满足绝对收敛求和表示（附录式 \eqref{eq:app-horizon-moment-sum}）：
\begin{equation}\label{eq:xi-horizon-moment-sum-recall}
\boxed{\;
c_n=\int_{\TT}\xi^{-n}\,\dd\mu_{\zeta}(\xi)
=
\sum_{\gamma\in\Gamma_+}\frac{\mathrm{mult}(\gamma)}{1+\gamma^2}\bigl(\xi_{\gamma}^{-n}+\xi_{\gamma}^{n}\bigr),
\qquad n\ge 0.
\;}
\end{equation}
其中 \(\Gamma_+\subset(0,\infty)\) 为 \(E_{\zeta}\) 的正零点多重集（按重数计）。
\par
权重 \((1+\gamma^2)^{-1}\) 的出现并非任选：它由 Cayley 坐标在边界上的 Jacobian（等价地，由对数导数极点在 Cayley 推回下的留数变换）结构性地强制给出；其语义后果是将高处零点在视界谱中压缩为弱耦合原子，并把可见性困难集中到端点邻域（参见附录定理 \ref{thm:app-horizon-endpoint-window-mass} 与推论 \ref{cor:app-offcritical-horocycle-busemann}）。

\paragraph{共动视界层析：平移族的洛伦兹窗离散谱}
为对“视界毛发”进行位置参数化扫描，引入共动平移族：对任意 \(x_0\in\RR\) 定义
\begin{equation}\label{eq:xi-horizon-comoving-carath-def}
\boxed{\;
\mathscr{C}_{\zeta,x_0}(w):=-\mathrm{i}\,m_{\zeta}\!\bigl(x_0+\mathcal{C}^{-1}(w)\bigr),
\qquad w\in\mathbb{D}.
\;}
\end{equation}
在 RH 下，\(m_{\zeta}\) 为 Herglotz 函数（附录定理 \ref{thm:app-horizon-weyl-herglotz}），而平移 \(z\mapsto z+x_0\) 为 \(\mathbb{H}\) 的自同构，故 \(\mathscr{C}_{\zeta,x_0}\) 对每个 \(x_0\) 仍为 Carath\'eodory 函数。
其谱测度可取为如下洛伦兹窗加权的纯点测度。

\begin{theorem}[共动视界层析的离散谱测度与矩恒等式]\label{thm:xi-horizon-comoving-spectral-measure}
假设 RH 成立。
令 \(\Gamma\subset\RR\) 为 \(E_{\zeta}\) 的全部实零点多重集（按重数计），并记其重数为 \(\mathrm{mult}(\gamma)\)。
对任意 \(x_0\in\RR\) 定义原子点
\[
  \xi_{\gamma,x_0}:=\mathcal{C}(\gamma-x_0)\in\TT\qquad(\gamma\in\Gamma).
\]
则 \(\mathscr{C}_{\zeta,x_0}\) 的 Carath\'eodory 表示可取为纯原子测度
\begin{equation}\label{eq:xi-horizon-mu-comoving}
\boxed{\;
\mu_{\zeta,x_0}
:=
\sum_{\gamma\in\Gamma}\frac{\mathrm{mult}(\gamma)}{1+(\gamma-x_0)^2}\,\delta_{\xi_{\gamma,x_0}}.
\;}
\end{equation}
并且其矩（视界毛发多极矩）满足
\begin{equation}\label{eq:xi-horizon-moment-comoving}
\boxed{\;
c_n(x_0):=\int_{\TT}\xi^{-n}\,\dd\mu_{\zeta,x_0}(\xi)
=\sum_{\gamma\in\Gamma}\frac{\mathrm{mult}(\gamma)}{1+(\gamma-x_0)^2}\,\xi_{\gamma,x_0}^{-n},
\qquad n\ge 0.
\;}
\end{equation}
特别地，
\begin{equation}\label{eq:xi-horizon-c0-im-m}
\boxed{\;
c_0(x_0)=\mu_{\zeta,x_0}(\TT)
=\sum_{\gamma\in\Gamma}\frac{\mathrm{mult}(\gamma)}{1+(\gamma-x_0)^2}
=\Im\,m_{\zeta}(x_0+\mathrm{i}).\;}
\end{equation}
\end{theorem}
\noindent
其推导见附录段落 \ref{par:app-horizon-comoving-spectral-measure}。

\begin{remark}[洛伦兹窗平滑与层析语义]\label{rem:xi-horizon-lorentz-tomography}
式 \eqref{eq:xi-horizon-c0-im-m} 表明：\(x_0\mapsto c_0(x_0)\) 是零点计数测度在核 \((1+t^2)^{-1}\) 下的卷积型平滑；
因此 \(c_0(x_0)\) 在 \(x_0\)-方向呈现以各个 \(\gamma\) 为中心的洛伦兹峰，并可作为“视界毛发总质量”的层析扫描信号。
该信号与附录定理 \ref{thm:app-horizon-endpoint-window-mass} 的端点角窗质量同源：二者均由同一 Cauchy/Poisson 权重 \((1+t^2)^{-1}\) 诱导，只是观测在 \(\TT\) 的端点角窗与 \(\RR\) 的平移窗之间切换。
\end{remark}

\begin{remark}[与端点压缩律的对照：共动局部化解除二次径向抑制]\label{rem:xi-horizon-comoving-vs-redshift}
若考虑离临界线偏移 \(\rho=\beta+\mathrm{i}\gamma\)（\(\delta=\tfrac12-\beta>0\)）在圆盘紧化下的端点层余量，则其固定图表深度满足二次压缩律
\(
h(\rho)=\frac{4\delta}{\gamma^2+(1+\delta)^2}
\)
（定理 \ref{thm:xi-offcritical-quadratic-radial-compression}）。
而在移位 Cayley 的共动图表 \(x_0=\gamma\) 下，深度退化为
\(
h_{\gamma}(\rho)=\frac{4\delta}{(1+\delta)^2}\sim 4\delta
\)
（附录定理 \ref{thm:app-comoving-horizon-localization}），从而偏移量以线性尺度显现而不再受 \(\gamma^{-2}\) 红移抑制。
该对照为“共动层析”的制度含义提供了解析刻画：共动参数把端点边界层中的局部结构重新对齐到可见尺度，但其代价与可寻址性条款见注 \ref{rem:xi-comoving-horizon-localization-address-null}。
\end{remark}

\paragraph{酉谱重构接口：相位化闭包与视界矩序列}
在 unitary-slice 的相位化闭包口径中，Carath\'eodory 表示不仅给出正性证书语言，也给出一个与有限维 CMV/Schur 参数化严格同型的酉谱实现接口。

\begin{corollary}[视界谱测度的酉谱实现（与 Ramanujan 相位化一致）]\label{cor:xi-horizon-unitary-spectral-realization}
假设 RH 成立并令 \(\mu_{\zeta}\) 为 \eqref{eq:xi-horizon-mu-atomic-recall} 的视界谱测度。
则存在一酉算子 \(U_{\zeta}\) 及循环向量 \(\mathbf{e}_0\)（可取为谱模型 \(L^2(\TT,\dd\mu_{\zeta})\) 上的常数函数 \(1\)），使得对一切 \(n\ge 0\) 有
\begin{equation}\label{eq:xi-horizon-unitary-moments}
\boxed{\;
c_n=\langle U_{\zeta}^{\,n}\mathbf{e}_0,\mathbf{e}_0\rangle. \;}
\end{equation}
同理，对任意 \(x_0\in\RR\) 与 \(\mu_{\zeta,x_0}\)（定理 \ref{thm:xi-horizon-comoving-spectral-measure}），亦存在对应的酉谱模型实现 \(c_n(x_0)=\langle U_{\zeta,x_0}^{\,n}\mathbf{e}_0,\mathbf{e}_0\rangle\)。
该谱化接口与文内有限维情形的“酉性 \(\Leftrightarrow\) 单位圆根”CMV 实现（命题 \ref{prop:xi-hilbert-polya-cmv}）同型：差异仅在于此处的极限对象为真实 \(\Xi_{\zeta}\) 所诱导的（可能）无限维谱测度。
\end{corollary}

\paragraph{与素数侧绝对收敛补丁拼接：加权显式公式族}
附录定理 \ref{thm:app-horizon-euler-patch} 给出：在 \(|w|<1/3\) 内，\(\mathscr{C}_{\zeta}(w)\) 可由伽马项与 \(\zeta'/\zeta\) 的 Dirichlet 级数以一致绝对收敛形式表示；推论 \ref{cor:app-horizon-prime-moment-interface} 因而给出 \(c_n\) 的可审计素数权系数接口。
与零点侧的离散谱公式 \eqref{eq:xi-horizon-moment-sum-recall} 并置，即得到一族“素数侧（绝对收敛）＝零点侧（红移权重相位和）”的严格恒等式。

\begin{corollary}[加权显式公式（视界矩版本）]\label{cor:xi-horizon-weighted-explicit-formula}
假设 RH 成立。
则对每个 \(n\ge 0\)，视界多极矩 \(c_n\) 同时满足：\((i)\) 由附录推论 \ref{cor:app-horizon-prime-moment-interface} 给出的绝对收敛素数权表达；\((ii)\) 由零点原子谱给出的红移加权相位和 \eqref{eq:xi-horizon-moment-sum-recall}。
此外，对每个 \(x_0\in\RR\)，共动矩 \(c_n(x_0)\) 亦具有同型的素数侧绝对收敛表达；其算术项等价于在 Dirichlet 级数中引入相位因子 \(n^{\mathrm{i}x_0}\)（即 \(n^{-(\sigma-\mathrm{i}x_0)}=n^{\mathrm{i}x_0}n^{-\sigma}\)），并与零点侧层析公式 \eqref{eq:xi-horizon-moment-comoving} 严格相等。
\end{corollary}
\noindent
共动版本在 \(|w|<1/3\) 内的绝对收敛欧拉积补丁见附录注 \ref{rem:app-horizon-euler-patch-comoving}。

\endinput
